\documentclass[10pt,a4paper]{amsart}
\usepackage[margin=19mm]{geometry}
\usepackage{amsmath,amssymb,mathtools}
\usepackage[scr=rsfs,cal=euler]{mathalfa}
\usepackage{xfrac}
\usepackage[unicode,hidelinks,bookmarks=false]{hyperref}
\newcommand{\ie}{i.e.\ }
\newcommand{\cf}{cf.\ }
\newcommand{\resp}{respectively\ }
\newcommand{\Subsection}{\subsection}
\providecommand{\nobreakdash}{\nobreak\hbox{-}}
\setlength{\emergencystretch}{2em}
\multlinegap0pt
\usepackage{enumerate}
\usepackage{amssymb}
\usepackage{csquotes}
\newtheorem{cor}{Corollary}
\title{A Class of algebras admitting infinitely many norm topologies}
\author{J. G. Patel}
\date{Source: arXiv:2602.16404v1 (CC BY 4.0)}
\begin{document}
\maketitle

\noindent\emph{Source and licence.} A Class of algebras admitting infinitely many norm topologies. Version arXiv:2602.16404v1 (CC BY 4.0), distributed by arXiv under the Creative Commons Attribution 4.0 International licence, http://creativecommons.org/licenses/by/4.0/, as recorded in the arXiv licence metadata for this exact version and shown on the abstract page https://arxiv.org/abs/2602.16404v1. Full licence notice: \texttt{licenses/arxiv-2602.16404-CC-BY-4.0.txt}.\par\smallskip

\noindent\textit{Editorial scope.} Counted unit: the article's cor Norms (source0.tex lines 126--128). The article's complete original proof of it is delivered with it (source0.tex lines 130--149, one \texttt{proof} environment of 20 lines). No mathematics is written, repaired or re-derived: the statement and the proof are the article's own lines, in the article's own notation, and no retained line is substituted or re-flowed.  No further source-local context is needed: every cross-reference of every retained line resolves inside the two delivered spans.  Nothing was omitted from any retained span, so the package carries no omission record.  No retained line contains a citation, so the wrapper carries no bibliography and no bibliography entry.  No retained line contains a figure environment, an image-inclusion command or drawing code, so no image is delivered and the package declares no asset dependency.  Editorial formatting only, in addition: an AMS \texttt{amsart} preamble that restates the article's own declarations and macros used by the retained lines, namely the environments \texttt{cor} on separate counters, together with the standard packages those lines need.  The retained environments print in this document as Corollary 1 in order of appearance, whereas the article prints the same items with the numbering of its own full document.  The complete native source of the article as distributed in the arXiv e-print for this exact version is bundled unchanged as \texttt{sources/194878/source0.tex}.
 \begin{cor}{\label{Infinitely Many Norms}}
 	Let $\mathcal{A}$ be a normed algebra, and let the codimension of $\mathcal{A}^2$ in  $\mathcal{A}$ is infinite. Then $\mathcal{A}$ admits infinitely many algebra norm.
 \end{cor}

 \begin{proof} Let $(\mathcal{A}, \| \cdot \|)$ be a normed algebra. Since $\mathcal{A}^2$ has infinite co-dimension in $\mathcal{A}$, there exists a countably infinite linearly independent subset $L$ of $\mathcal{A}$ such that $\mathcal{A}^2 \cap L = \phi$ and $\| a \| =1 \; (a \in L)$. For each $n \in \mathbb{N}$, choose $L_n = \{ a_{n1}, a_{n2}, \ldots \} \subset L$ such that: 
 	\begin{enumerate}
 		\item Each $L_n$ is infinite;
 		\item $L_n \cap L_m = \phi \ (n \neq m)$;
 		\item  $L = \bigcup_{n=1}^{ \infty} L_n$.
 	\end{enumerate}
 	Let $D$ be a basis of $\mathcal{A}^2$. Let $B_n$ be a (Hamel) basis of $\mathcal{A}$ such that $L_n \cup D \subset B_n$ for each $n \in \mathbb{N}$. Now, consider the (unique) linear map $\varphi_n : \mathcal{A} \longrightarrow \mathbb{C}$ such that
 	\[
 	\varphi_n(a)=
 	\begin{cases}
 		k&  (\text{if } a=a_{nk}  \in L_n \setminus \{a_{n1}\});\\
 		1 & (\text{if } a \in B_n \setminus (L_n \cup D)); \\
 		0 & (\text{if } a \in D \cup \{a_{n1}\}).
 	\end{cases}
 	\]
 	Now for $a \in \mathcal{A}$, define
 	$$p_n(a)=\|a\|+ |\varphi_n(a)|.$$
 	Then $(\mathcal{A}, p_n( \cdot ))$ is a normed algebra.
 	Clearly, each $p_n(\cdot)$ is a linear norm. Let $a, b \in \mathcal{A}$. Then $p_n(ab)  = \|ab\| \leq p_n((a, \alpha))p_n((b, \beta))$ because $ab \in \mathcal{A}^2$ and hence $\varphi_n(ab) = 0$. Thus each $p_n(\cdot) \in N(\mathcal{B})$. Now, we \textit{claim} that these norms are non-equivalent on $\mathcal{A}$. Let $m < n$ and $g_k = a_{mk} \ (k \in \mathbb{N})$. Then, for $k \geq 2, \; p_m((a_k, 0)) = 1+k$ and $p_n((a_k, 0)) \leq 2$ because $\varphi_n(a_{mk}) = 0 \text{ or 1}$. Thus we have proved our claim.
 \end{proof}

\end{document}
